% \cventry{<name> \cvlink{<url>}{Code}}{<year>}{<one-line subtitle, optional>}{}
% \begin{cvbullets} \item ... \end{cvbullets}
\cvsection{Software Development}

\cventry{Learn-MPM-2D\cvlinksep\cvlink{https://yjchoi1.github.io/learn-mpm-2d/main_page.html}{Website}}{2024}{}{}
\begin{cvbullets}
  \item Developed an open-source learning package for the 2-D material point method, guiding learners from core theory to hands-on practice in a beginner-friendly interactive Python notebook environment.
\end{cvbullets}

\cventry{Cbgeopy: Material Point Method Model Generator\cvlinksep\cvlink{https://cbgeopy.readthedocs.io/en/latest/}{Docs}}{2024}{}{}
\begin{cvbullets}
  \item Developed a material point method model generator that supports various geometries from simple columns to multi-layered 3-D topography, which can be integrated with CB-Geo MPM.
\end{cvbullets}

\cventry{GNS: A Generalizable Graph Neural Network-Based Simulator for Particulate and Fluid Modeling\cvlinksep\cvlink{https://github.com/geoelements/gns}{Code}}{2023}{}{}
\begin{cvbullets}
  \item Contributed as a core developer to implement mesh graph networks for fluid modeling, multi-material features, and code maintenance.
\end{cvbullets}

% \cventry{P-Wave Arrival Time Detector, DesignSafe Academy/NHERI Hackathon\cvlinksep\cvlink{https://github.com/yjchoi1/pwave_dancers}{Code}}{2022}
%   {Texas Advanced Computing Center}{USA}
% \begin{cvbullets}
%   \item Developed a p-wave arrival time detector for earthquake sources using LSTM and CNN models as a machine learning model developer in a team hackathon on machine learning and high-performance computing.
% \end{cvbullets}
